\chapter{53}



Im kleinen Schrank neben dem Herd fand Louis eine Auswahl an Teebeuteln.
Er setzte Wasser auf. 

Mit der Tasse in der Hand stellte er sich in den
Türrahmen der Hütte. Die Sonne wirkte nur mehr als Ahnung nach. Wie eine gute Freundin, die sich
verabschiedet hat. Mit ihr waren Licht und Wärme so
selbstverständlich verschwunden, wie sie in einigen Stunden wieder auftauchen würden.
Eigentlich passte doch gerade alles. Wenn er nur dieses negative
Grundflattern ausschalten könnte. Louis’ Körper spannte sich an. Die
altbekannte Traurigkeit war wieder da. Als verteidige sie als
Zurückgestossene ihren Platz. Und all seine Vorhaben für den Abend
schmolzen urplötzlich als Nichtigkeiten dahin. Ihr würde diese Einsamkeit folgen,
die er nur zu gut kannte.

«Tatata-tatata-tata-ta». Die laute Marimba-Melodie riss ihn aus seinen
Gedanken. Erschrocken griff er nach seinem Handy und starrte auf das
Display. Mazi. Natürlich. Er fasste sich an die Stirn. Es galt, ihr
Treffen von morgen zu vereinbaren.

«Mazi, hallo.»

Louis richtete sich auf.

«Hi, Ciro, gut angekommen?»

Mazi klang nah. Als stünde er direkt neben ihm.

«Ja, danke. Alles okay. Ich bin schon fast bereit für die Nacht. Nur die
Dusche hab ich noch nicht eingeweiht,» Louis versuchte, unbekümmert zu
klingen, «hab auch an dich gedacht, wie machen wir’s? Soll ich morgen zu
dir rüber kommen oder kommst du zum Frühstück zu mir?»

«Ich glaub, es ist einfacher, wenn ich zu dir komme. Zum Frühstück,
klingt gut, danke. Um sechs?»

Das war früh.

«Gute Idee. Ich bereite uns was vor.»

«Toll,» Mazi wirkte vergnügt, «ich wünsche dir eine entspannte, erste
Alpnacht, bis morgen, Ciro, ciao.»

«Danke, danke, bis morgen. Ciao.»

Louis trat in die Hütte.
Während er die Leiter in den Schlafbereich hochkletterte, dachte er an
Mazi. Ein sympathischer Kerl. Ihre Verabredung morgen war ein
willkommener Fixpunkt. Vielleicht sollte er seine Traurigkeit einfach in
Ruhe lassen.

Mit einem Frottiertuch über der Schulter trat er auf den kleinen Vorplatz. 
Er freute sich auf seine erste Dusche unter freiem Himmel.
Der Abend hatte sich unterdessen wie eine Glocke über die Umgebung
gestülpt. Damit wandelte sich die Alpstimmung erneut. Als er nackt unter der Brause stand, empfand er die Dunkelheit wie eine warme Umarmung. Der feine Duschstrahl war wohltuend.

Bevor er sich schlafen legte, steckte er sich
die Kopfhörer in die Ohren.
Die niedrige Zimmerhöhe liess die Decke wie ein Schutzschild von oben wirken. Er
sog den Duft des Raumes in sich auf. Dann begann er, seine 
Erinnerung nach passender Musik abzusuchen. Er wollte sich beruhigt in
den Schlaf hören. Morgen musste er früh raus. Das Handy in der Hand schloss er die Augen. Mit jedem
Atemzug fühlte er sich schwerer. Unverhofft tauchte das Bild auf. Er in
Grandpères Arm, geborgen in alter Zeit. Die Wärme in der Wohnung, hin
und wieder Geräusche aus der Küche. Maman, die gelegentlich zufrieden ins
Wohnzimmer schaute, Lucia am Boden, 
vor sich den hölzernen Bauernhof. Und Stings Musik, die er lieben
gelernt hatte. Die Texte verstand er nicht. Es reichte, dass Grandpère
die Songs mochte und Louis sich dabei wohl fühlte. Stings
\emph{«Fragile»} hatte ihm von Anfang an am besten gefallen.
Louis gab den Titel ein. Augenblicklich fühlte er sich gut. Wie
damals.
Von weit weg erreichte ihn Stings Stimme, passend zu den gedämpften
Klängen. Er schloss die Augen. Weitere Bilder kamen. Er sah sich am Boden vor dem Lautsprecher sitzen.
Er konnte höchstens acht gewesen sein. Er hatte seine Gitarre geholt und versucht, das Stück nachzuspielen. Hatte es nicht geschafft und das Instrument enttäuscht in die Ecke gestellt. Erst ein gutes Jahrzehnt später spielte er den Song, nicht überragend gut, doch ganz  passabel. Bevor er allmählich einschlummerte, lächelte er.

\qrblock{Sting \emph{«Fragile»}}{media/QR-Codes/037_Fragile_Sting_Spotify.png}  
